% Independent tau=-1 degree-three ramification exclusion.
% Prepared 2026-08-02 against the Program 2 (3,4) Hilbert--Burch chart.

\begin{theorem}[Primitive $\tau=-1$ chart]
\label{thm:quartic-r3-tau-minus-one}
Let $k$ be algebraically closed of characteristic zero, and let
\[
 F=LX+H_2+H_3+H_4:\mathbb A_k^3\longrightarrow\mathbb A_k^3
\]
be Keller.  Assume the preceding Program~2 reductions place $F$ in the
primitive coprime binary-pencil branch of common ramification degree three,
and in the independent Hilbert--Burch chart
\[
 u=x^2,\qquad v=y^2,\qquad
 w=ax^3+x^2y+xy^2+by^3,\qquad \tau=-1.
\]
Then the first normal layer vanishes.  Consequently this chart exits to the
zero-normal triangular branch and contains no Keller counterexample.
\end{theorem}

\begin{proof}
The differential Hilbert--Burch equation
\[
 x^2\,dP-y^2\,dQ+w\,dR=0
\]
reconstructs
\begin{align*}
 R={}&-ax^3-3x^2y+3xy^2+by^3,\\
 P=\frac34\bigl(&a^2x^4-2abxy^3+4ax^3y-4ax^2y^2
                    -2by^4+2x^2y^2-2xy^3-y^4\bigr),\\
 Q=-\frac34\bigl(&2abx^3y+2ax^4-b^2y^4+4bx^2y^2-4bxy^3
                    +x^4+2x^3y-2x^2y^2\bigr).
\end{align*}
Their minors have the required common cubic:
\[
 J(Q,R)=x^2\Delta,\qquad J(P,R)=y^2\Delta,
 \qquad J(P,Q)=w\Delta.
\]
Put
\[
 F_1=27a^2b^2-18ab+4a+4b-1,
 \qquad
 G=a^2b^2-6ab-4a-4b-3.
\]
Exact dehomogenized resultants give
\[
 \operatorname{Res}_x(P(x,1),Q(x,1))
   =-\frac{6561}{65536}F_1G^3.
\]
Thus primitivity implies $F_1G\ne0$.

Write $A=(H_3)_1$, $B=(H_3)_2$, and $E=(H_2)_3$.  The first normal
syzygy has the complete form
\begin{equation}
\label{eq:tau-minus-one-first-normal}
 (A_z,B_z,E_z)
 =L_0(-ax-y,x+by,1)+\mu(y^2,x^2,0),
 \qquad L_0=l_0x+l_1y+l_2z.
\end{equation}
Let $D_j$ denote the $j$th coefficient of the determinant arc, using
$H_4,H_3,H_2,H_1$ as levels $0,1,2,3$.  The coefficient of $z^3$ in
$D_3$ is
\[
 \frac{3l_2^2}{4}\bigl(
 A_*x^3+C_*x^2y-D_*xy^2-B_*y^3\bigr),
\]
where
\[
 A_*=a^2b-3a-2,\quad B_*=ab^2-3b-2,
 \quad C_*=ab+2a+1,\quad D_*=ab+2b+1.
\]
If $l_2\ne0$, all four coefficients vanish.  Since
\[
 C_*-D_*=2(a-b),\qquad C_*\big|_{b=a}=(a+1)^2,
\]
this forces $(a,b)=(-1,-1)$, where $G=0$.  Hence $l_2=0$ on the
primitive locus.

Set $\lambda=l_0x+l_1y$.  The coefficient of $z$ in $D_3$ gives six
quadrics in
\[
 l_0^2,\ l_0l_1,\ l_1^2,\ l_0\mu,\ l_1\mu,\mu^2,\ u_2,\ v_2,
\]
where $u_2,v_2$ are the $z^2$ coefficients of the first two components
of $H_2$.  The first and last equations are
\[
 A_*(l_0^2-6u_2)=0,\qquad
 B_*(l_1^2-6v_2)=0.
\]
We divide the remaining analysis into the complete coefficient-divisor
stratification.

\smallskip
\noindent\emph{The open set $A_*B_*\ne0$.}
Here $u_2=l_0^2/6$ and $v_2=l_1^2/6$.  After removing harmless scalar
factors, the four middle equations generate the ideal $I=(f_0,f_1,f_2,f_3)$,
where
\begin{align*}
 f_0={}&a^2bl_0l_1-abl_0^2-2al_0^2-3al_0l_1+4al_0\mu-2a\mu^2
        -l_0^2-2l_0l_1+4l_0\mu,\\
 f_1={}&a^2bl_1^2-2abl_0^2+abl_0l_1+4abl_0\mu+2al_0l_1-3al_1^2
        +4al_1\mu-4bl_0^2\\
      &\hspace{1.4em}-2l_0^2+l_0l_1-4l_0\mu-2l_1^2+4l_1\mu-6\mu^2,\\
 f_2={}&ab^2l_0^2+abl_0l_1-2abl_1^2-4abl_1\mu-4al_1^2-3bl_0^2
        +2bl_0l_1-4bl_0\mu\\
      &\hspace{1.4em}-2l_0^2+l_0l_1-4l_0\mu-2l_1^2+4l_1\mu-6\mu^2,\\
 f_3={}&ab^2l_0l_1-abl_1^2-3bl_0l_1-2bl_1^2-4bl_1\mu-2b\mu^2
        -2l_0l_1-l_1^2-4l_1\mu.
\end{align*}
Exact Gröbner reduction gives the saturation certificates
\begin{equation}
\label{eq:tau-minus-one-saturation}
 \mu^4G\in I,
 \qquad
 l_0^4G,\ l_1^4G\in I+(\mu).
\end{equation}
Thus every nonzero projective point $[l_0:l_1:\mu]$ in this open set
forces $G=0$, contrary to primitivity.

\smallskip
\noindent\emph{The divisor $A_*=0$, with $B_*\ne0$.}
Here $a\ne0$ and
\[
 b=\frac{3a+2}{a^2},\qquad
 B_*=-\frac{2(a-2)(a+1)^2}{a^3}.
\]
Use $v_2=l_1^2/6$ in the four middle equations.  On the three standard
projective charts $l_0=1$, $l_1=1$, and $\mu=1$, their exact ideals contain,
respectively,
\[
 (a+1)^4,\qquad a^2(a+1)^4,\qquad a^2.
\]
Since $a\ne0$ and $B_*\ne0$, none has a point.  The divisor $B_*=0$,
$A_*\ne0$, is checked independently: after
$a=(3b+2)/b^2$ and $u_2=l_0^2/6$, the corresponding three chart ideals
contain
\[
 b^2(b+1)^4,\qquad (b+1)^4,\qquad b^2.
\]
It likewise has no point.

\smallskip
\noindent\emph{The intersection $A_*=B_*=0$.}
The identity
\[
 A_*-B_*=(a-b)(ab-3)
\]
shows that $ab=3$ would give $A_*=-2$, so $a=b$.  Therefore
\[
 (a,b)=(-1,-1)\quad\hbox{or}\quad (2,2).
\]
The first point has $G=0$.  At the primitive point $(2,2)$, elimination of
$u_2,v_2$ leaves
\begin{align*}
 K_1&=9l_0^2-9l_0l_1-12l_1\mu+2\mu^2=0,\\
 K_2&=9l_0l_1-12l_0\mu-9l_1^2-2\mu^2=0.
\end{align*}
These satisfy
\[
 K_1-K_2=(3l_0-3l_1+2\mu)^2,
 \qquad
 K_1+K_2=3(l_0+l_1)(3l_0-3l_1-4\mu).
\]
Consequently the nonzero first-normal solutions are exactly
\begin{align*}
 (l_0,l_1,\mu,u_2,v_2)
   &=\left(h,h,0,\frac{h^2}{18},\frac{h^2}{18}\right),\\
 (l_0,l_1,\mu,u_2,v_2)
   &=\left(-h,h,3h,-\frac{5h^2}{6},-\frac{5h^2}{6}\right),
 \qquad h\ne0.
\end{align*}

It remains to test the next determinant coefficient.  Keep arbitrary binary
cubics $A_0,B_0$, arbitrary binary quadrics $C_0,D_0,E_0$, arbitrary binary
linear forms $p,q$, and an arbitrary linear part.  Thus
\[
 A=A_0+zA_z,\quad B=B_0+zB_z,
\]
\[
 H_2=\bigl(C_0+zp+u_2z^2,
           D_0+zq+v_2z^2,
           E_0+z\lambda\bigr).
\]
A direct exact expansion, with no specialization of these lower terms, gives
for the two displayed branches
\begin{align*}
 [z^2]D_4
   &=-\frac{h^3}{2}(x+y)(5x^2+2xy+5y^2),\\
 [z^2]D_4
   &=\frac{135h^3}{2}(x-y)(x+y)^2.
\end{align*}
Both are nonzero for $h\ne0$.  Hence the intersection has no nonzero
first-normal solution either.

We have proved that primitivity forces
$l_0=l_1=l_2=\mu=0$.  The established zero-normal lemma then makes the third
coordinate triangular (or makes every nonlinear term binary), and the usual
plane Keller reduction proves that $F$ is an automorphism.
\end{proof}

\begin{remark}[Computation boundary]
The accompanying checker reconstructs $P,Q$, the six first-normal equations,
the saturation and projective-chart certificates, and the two unrestricted
$[z^2]D_4$ obstructions over $\mathbb Q$.  It does not re-prove the upstream
placement of a quartic map in this Hilbert--Burch chart or the final
zero-normal plane theorem.
\end{remark}
